\section{Introduction}
\label{sec:intro}

Cell-free DNA (cfDNA) whole-genome sequencing has become the dominant non-invasive substrate for non-invasive prenatal testing (NIPT), minimal residual disease (MRD) monitoring, and early cancer detection \citep{snyder2016,sun2015,moss2018,loyfer2023}. The clinical pipeline produces enormous quantities of low-pass WGS data, but the vast majority of methylation information remains inaccessible because bisulfite conversion is destructive, expensive, and biased toward fragment ends. \citet{liu2024finaleme} closed part of this gap with FinaleMe by showing that fragmentomic features---fragment length, end motifs, GC content, and orientation---are sufficient to recover per-CpG methylation status through a per-sample hidden Markov model. FinaleMe demonstrates genome-wide Pearson correlations that vary by region between roughly $0.5$ in CpG-sparse regions and approximately $0.7$ in CpG-rich regions on $20\times$ low-pass WGS, with substantial degradation at lower coverage and outside CpG islands; we refer the reader to \citet[Fig.~3]{liu2024finaleme} for the full empirical breakdown.

Despite its impact, FinaleMe has well-defined limitations that motivate this work. We list them as seven distinct items, each tied to the corresponding HAVI-Methyl design choice that addresses it.

\paragraph{(L1) Binary latent state.} FinaleMe's hidden state is binary $\{m,u\}$, discarding the continuous $\beta\in[0,1]$ structure that drives all downstream tissue-of-origin and biological-age inference. \emph{HAVI-Methyl response:} a continuous logit-$\beta$ latent with a Beta-Binomial likelihood (Section~\ref{sec:model}).

\paragraph{(L2) No uncertainty quantification.} FinaleMe returns Viterbi point estimates, which is a serious obstacle for clinical deployment where calibrated risk is required. \emph{Response:} a normalizing-flow variational posterior with conformal coverage guarantees (Sections~\ref{sec:varfamily} and~\ref{sec:calibration}).

\paragraph{(L3) Per-sample fitting.} FinaleMe re-runs Baum-Welch from scratch on every sample, ignoring statistical strength shared across thousands of samples in modern biobanks. \emph{Response:} a hierarchical Bayesian model with a population locus prior, fitted by stochastic variational inference (Sections~\ref{sec:model} and~\ref{sec:elbo}).

\paragraph{(L4) Buffy-coat prior leakage.} The prior is seeded from the same individual's leukocyte methylome, so any disease signal that also alters leukocytes (inflammation, treatment effects) is absorbed into the prior and erased from the prediction. \emph{Response:} a variational information-bottleneck penalty, an mQTL-anchor instrumental-variable loss, and counterfactual augmentation (Section~\ref{sec:identify}).

\paragraph{(L5) Poor performance in CpG-poor regions.} The Markov chain along loci forces information to propagate through chromosomal neighbors that themselves lack signal. \emph{Response:} amortized inference over the fragment bag at each locus through a permutation-invariant Set Transformer encoder and a long-range DNA sequence context (Section~\ref{sec:varfamily}).

\paragraph{(L6) Tissue deconvolution on binarized predictions.} FinaleMe's downstream tissue-of-origin pipeline binarizes the per-CpG predictions before solving a quadratic program for tissue fractions, discarding posterior mass and breaking calibration. \emph{Response:} a Dirichlet head on continuous posterior means with optional hierarchical Dirichlet process extension, trained jointly with the methylation ELBO (Section~\ref{sec:too}).

\paragraph{(L7) No inter-fragment pooling.} Beyond the per-fragment HMM emissions, no mechanism aggregates information across fragments at a single locus. \emph{Response:} the Set Transformer pooling step that produces a permutation-invariant locus embedding from the entire fragment bag at once (Section~\ref{sec:varfamily}).

\paragraph{Contributions.} The paper makes the following contributions.

\noindent (i) A complete generative model for cfDNA fragmentomics with a hierarchical logit-Normal prior and Beta-Binomial / Categorical / Negative-Binomial likelihoods (Section~\ref{sec:model}).

\noindent (ii) A full ELBO derivation with approximate natural-gradient updates on the population latents and an amortized neural-spline-flow posterior on the per-(sample, locus) latents, with explicit handling of the non-conjugate natural-gradient correction (Sections~\ref{sec:elbo} and~\ref{sec:algo}).

\noindent (iii) Five formal propositions and two new theoretical contributions, including hierarchical identifiability under instrumental anchors with explicit nonparametric IV completeness, an asymptotic variational-information-bottleneck prior-leakage bound with a finite-$\beta_{\mathrm{VIB}}$ corollary, conformal marginal coverage, convergence of non-conjugate natural-gradient SVI, a hierarchical-pooling variance-reduction inequality, and a Fano-style information-theoretic recoverability lower bound (Section~\ref{sec:theory}).

\noindent (iv) A chromatin-aware cfDNA fragmentation simulator that generates paired ground-truth methylation and WGS reads (Section~\ref{sec:simulator} and Appendix~\ref{app:sim}).

\noindent (v) Synthetic recovery experiments with reproducible code (Section~\ref{sec:synth}).

\noindent (vi) A benchmarking plan against FinaleMe on the 80 paired WGS/WGBS samples and the Sun, Moss, and Loyfer reference atlases, with full ablations (Section~\ref{sec:benchmark}).

\noindent (vii) A wafer-scale accelerator training recipe with parameter-count and FLOP estimates exploiting wafer-scale weight streaming for genome-wide locus mini-batching (Section~\ref{sec:benchmark}).

\paragraph{Relationship to methylation foundation models.} HAVI-Methyl complements rather than competes with the recent foundation-model line of work on methylation, of which CpGPT \citep{cpgpt2024} and MethylGPT \citep{methylgpt2024} are the leading examples. Foundation models consume per-CpG $\beta$-values from methylation arrays as input and produce embeddings or downstream predictions; HAVI-Methyl consumes raw cfDNA WGS fragments and produces $\beta$-values as output. The two pipelines compose: HAVI-Methyl provides the input to a foundation model, and the foundation model provides downstream embeddings, biological-age clocks (PhenoAge, GrimAge, DunedinPACE, PC-clocks; \citealp{horvath2013,levine2018phenoage,lu2019grimage,belsky2022dunedinpace,higginschen2022pcclocks}), or imputation of missing CpGs. Section~\ref{sec:related} expands on this relationship.

\paragraph{Roadmap.} Section~\ref{sec:background} reviews cfDNA fragmentomics, variational Bayesian inference, and the methylation-clock and foundation-model literature. Section~\ref{sec:model} specifies the generative model. Section~\ref{sec:varfamily} defines the structured variational family. Section~\ref{sec:elbo} derives the ELBO with all reparameterization, IWAE, mini-batch rescaling, and natural-gradient details. Section~\ref{sec:algo} states the SVI training algorithm. Section~\ref{sec:identify} establishes identifiability through the VIB penalty, mQTL anchors, counterfactual augmentation, and domain-adversarial training. Section~\ref{sec:calibration} derives the conformal calibration wrapper. Section~\ref{sec:too} attaches the joint tissue-of-origin head. Section~\ref{sec:simulator} introduces the cfDNA simulator. Section~\ref{sec:synth} reports synthetic recovery results. Section~\ref{sec:benchmark} lays out the benchmarking plan on real data. Section~\ref{sec:theory} collects the seven formal propositions. Section~\ref{sec:discussion} discusses limitations and extensions. Section~\ref{sec:related} treats related work in detail. Section~\ref{sec:conc} concludes. Appendices~\ref{app:elbo}--\ref{app:extres} contain full derivations, the VB-HMM baseline updates, simulator validation, hyperparameter tables, the notation glossary, and extended results.
